%! Solving Radical Equations \& Extraneous Solutions — Unit 6, Lesson 6.5 — Guided Notes
\documentclass[10pt]{article}
\usepackage{algebra2-article}
\usepackage{algebra2-boxes}
\newcommand{\TallMath}[1]{$\displaystyle #1\rule[-1.4em]{0pt}{3.2em}$}
\newcommand{\lgrid}[4]{%
  \draw[step=1, linegray!40, very thin] (#1,#3) grid (#2,#4);
  \draw[->, semithick, charcoal] (#1,0)--(#2,0) node[below right, font=\scriptsize]{$x$};
  \draw[->, semithick, charcoal] (0,#3)--(0,#4) node[left, font=\scriptsize]{$y$};}

\begin{document}

\pageheader{Unit 6, Lesson 6.5}{Guided Notes}
\namedateperiod

\vspace{0.08in}

\begin{objectivebox}
\textbf{By the end of this lesson, I will be able to\dots}
\begin{itemize}\setlength{\itemsep}{2pt}
  \item solve an equation containing one radical by \emph{isolating} the radical and raising both sides to the $n$th power, and solve an equation of the form $x^{m/n}=k$;
  \item \emph{check} every candidate in the original equation, identify the \textbf{extraneous} ones, and verify a solution set on a graph;
  \item \emph{justify} why an extraneous solution appears --- and why an equation with an odd index never produces one.
\end{itemize}
\end{objectivebox}

\vspace{0.05in}

\begin{vocabbox}
\small
Fill in each term as we build it together.
\par\vspace{2pt}   % \par is required: \termblanklong's \noindent is a no-op mid-paragraph
\termblanklong{Radical equation}
\termblanklong{Isolating the radical}
\termblanklong{Candidate (possible solution)}
\termblanklong{Extraneous solution}
\termblanklong{The check}
\termblanklong{Graphical solution}
\end{vocabbox}

\begin{hookbox}
Here is a student's complete solution to $\sqrt{x+7}=x-5$. \textbf{Every line is correct.} Check
them one at a time and you will not find a mistake anywhere.
\begin{center}
\begin{tcolorbox}[colback=white, colframe=charcoal!45, boxrule=0.7pt, arc=1.5mm, width=0.80\linewidth,
    left=3mm, right=3mm, top=1.5mm, bottom=1.5mm]
\centering\small
\renewcommand{\arraystretch}{1.4}
\begin{tabular}{@{}l l@{}}
$\sqrt{x+7}=x-5$ & the original equation \\
$x+7=(x-5)^{2}$ & square both sides \\
$x+7=x^{2}-10x+25$ & expand \\
$0=x^{2}-11x+18$ & collect \\
$0=(x-2)(x-9)$ & factor (Warm-Up item 3) \\
$x=2$ \; or \; $x=9$ & \\
\end{tabular}
\end{tcolorbox}
\end{center}
\vspace{2pt}
Now test both answers in the \emph{original} equation:
\begin{center}
\renewcommand{\arraystretch}{1.5}
\begin{tabularx}{0.92\linewidth}{@{}c X X c@{}}
\hline
 & \textbf{Left side, $\sqrt{x+7}$} & \textbf{Right side, $x-5$} & \textbf{Equal?} \\
\hline
$x=9$ & \blank{2.0cm} & \blank{2.0cm} & \blank{1.2cm} \\
$x=2$ & \blank{2.0cm} & \blank{2.0cm} & \blank{1.2cm} \\
\hline
\end{tabularx}
\end{center}
\begin{itemize}\setlength{\itemsep}{1pt}
  \item One answer is wrong, and yet no \emph{step} was wrong. So where did the wrong answer come from? \writeline
\end{itemize}
Today's rule, and it is not optional: \textbf{squaring hands you candidates, not solutions. The
check is part of the method.}
\end{hookbox}

\begin{notesbox}{1. The Method: Isolate, Raise, Solve, Check}
A \textbf{radical equation} is an equation with the variable \emph{underneath} a radical. To get at
that variable you have to undo the radical --- and the way to undo an $n$th root is to raise both
sides to the \blank{1.6cm} power, because $\left(\sqrt[n]{A}\,\right)^{n}=\blank{1.0cm}$
(Warm-Up item 1).
\begin{center}
\begin{tcolorbox}[colback=goldbg, colframe=goldacc, boxrule=0.8pt, arc=1.5mm, width=0.94\linewidth,
    left=3mm, right=3mm, top=2mm, bottom=2mm]
\small
\textbf{Step 1.} \blank{3.2cm} the radical --- get it alone on one side, with nothing added to it.\\[3pt]
\textbf{Step 2.} Raise \emph{both sides} to the power equal to the \blank{1.6cm}.\\[3pt]
\textbf{Step 3.} Solve the equation that is left (it will be linear or quadratic).\\[3pt]
\textbf{Step 4.} \textbf{Check every answer in the \blank{2.6cm} equation.} \emph{Not optional.}
\end{tcolorbox}
\end{center}

\vspace{3pt}
\textbf{Why Step 1 comes first.} Look again at Warm-Up item 1(c). If the radical is \emph{not}
alone, squaring does not remove it:
\[
  \left(\sqrt{x}+3\right)^{2}=\blank{3.2cm}\ \ \text{--- still has a radical in it.}
\]
Squaring only erases a radical that is standing by itself.

\vspace{4pt}
\textbf{Worked example.} $\sqrt{x+4}-1=3$ \quad (this is $f(x)=\sqrt{x+4}-1$ from Lesson 6.0 --- we
are asking \emph{where does it equal $3$?})
\begin{center}
\renewcommand{\arraystretch}{1.7}
\begin{tabular}{@{}l l@{}}
$\sqrt{x+4}-1=3$ & \\
$\sqrt{x+4}=\blank{1.4cm}$ & Step 1: add $1$ to both sides --- the radical is now alone \\
$x+4=\blank{1.4cm}$ & Step 2: square both sides \\
$x=\blank{1.4cm}$ & Step 3: solve \\
\end{tabular}
\end{center}
\textbf{Step 4 --- check.} \; $\sqrt{\blank{1.2cm}}-1=\blank{1.0cm}-1=\blank{1.0cm}$. \;
True, so $x=\blank{1.2cm}$ really is a solution. \; \checkmark
\end{notesbox}

\begin{notesbox}{2. Why the Check Exists: Squaring Is a One-Way Street}
Warm-Up item 2 already said it. Squaring an equation is a perfectly legal move, but it is
\textbf{not reversible}:
\begin{center}
\begin{tcolorbox}[colback=forestbg, colframe=forest, boxrule=0.7pt, arc=1.5mm, width=0.94\linewidth,
    left=3mm, right=3mm, top=1.5mm, bottom=1.5mm]
\centering\small
If $a=b$, then $a^{2}=b^{2}$ --- \emph{always}. \\[3pt]
But if $a^{2}=b^{2}$, then $a=b$ \quad \textbf{or} \quad $a=\blank{1.4cm}$.
\end{tcolorbox}
\end{center}
So squaring can never \emph{lose} a solution, but it can \blank{2.6cm} solutions that the original
equation never had. Those invented answers are called \textbf{extraneous solutions}.

\vspace{3pt}
\textbf{Where the sign information goes.} $\sqrt{9}$ and $-3$ are different numbers. But
$\left(\sqrt{9}\,\right)^{2}$ and $(-3)^{2}$ are the \emph{same} number, $9$. Squaring cannot tell
$3$ from $-3$ --- it \textbf{destroys sign information}, and an extraneous solution is what comes
back through that hole.

\vspace{4pt}
\textbf{The Hook, finished.} $\sqrt{x+7}=x-5$ gave the candidates $x=2$ and $x=9$. Here is the
picture of what happened.
\begin{center}
\begin{minipage}[c]{0.48\linewidth}
\centering
\begin{tikzpicture}[scale=0.345]
  \lgrid{-8}{13}{-6}{6}
  \begin{scope}
    \clip (-8,-6) rectangle (13,6);
    \draw[very thick, forest, domain=-7:13, samples=110, smooth] plot (\x,{sqrt((\x)+7)});
    \draw[very thick, navy, domain=-1:11, samples=2] plot (\x,{(\x)-5});
  \end{scope}
  \draw[dashed, redacc, thick] (2,3)--(2,-3);
  \fill[forest] (-7,0) circle (4pt);
  \fill[charcoal] (9,4) circle (3.4pt) node[above left, font=\scriptsize]{$(9,4)$};
  \fill[redacc] (2,3) circle (3pt) node[above left, font=\scriptsize]{$(2,3)$};
  \fill[redacc] (2,-3) circle (3pt) node[below right, font=\scriptsize]{$(2,-3)$};
  \node[forest, font=\scriptsize] at (-3.6,4.3) {$y=\sqrt{x+7}$};
  \node[navy, font=\scriptsize] at (8.6,-3.6) {$y=x-5$};
\end{tikzpicture}
\end{minipage}
\hfill
\begin{minipage}[c]{0.48\linewidth}
\small
The two graphs cross \blank{1.4cm} time(s), at $x=\blank{1.0cm}$. \\[4pt]
At $x=2$ the graphs are \emph{not} equal: the curve is at $\blank{1.0cm}$ and the line is at
$\blank{1.2cm}$. \\[5pt]
Those two heights are opposites --- and squaring turns both of them into
$\blank{1.0cm}$. That is exactly the moment $x=2$ sneaked in. \\[5pt]
$x=2$ is \blank{2.6cm}. \\[4pt]
Solution set: $\{\blank{1.0cm}\}$
\end{minipage}
\end{center}
\vspace{2pt}
\begin{center}
\begin{tcolorbox}[colback=redbg, colframe=redacc, boxrule=0.7pt, arc=1.5mm, width=0.94\linewidth,
    left=3mm, right=3mm, top=1.5mm, bottom=1.5mm]
\small \textbf{You have seen this before.} In Lesson 5.6 you multiplied both sides of a rational
equation by an expression that might be \emph{zero}, and got candidates that had to be thrown out.
Today you square both sides of an equation whose sides might be \emph{opposite in sign}, and get
candidates that have to be thrown out. Same structure, different trapdoor: \emph{a step that is not
reversible produces candidates, and only the check turns a candidate into a solution.}
\end{tcolorbox}
\end{center}
\end{notesbox}

\begin{notesbox}{3. A Shortcut You Can See Before You Write Anything}
From Lesson 6.1: $\sqrt{\phantom{x}}$ means the \textbf{principal} root, and a principal even root
is never \blank{2.0cm}. So an equation that asks a square root to \emph{equal} a negative number
has \blank{2.6cm} --- and you can say so without doing any algebra at all.
\begin{center}
\begin{minipage}[t]{0.46\linewidth}
\small
\textbf{Read it first:} \; $\sqrt{2x-1}+5=2$
\[ \sqrt{2x-1}=\blank{1.4cm} \]
A square root can never be that. \\[3pt]
\textbf{Answer:} \blank{3.0cm}
\end{minipage}
\hfill
\begin{minipage}[t]{0.50\linewidth}
\small
\textbf{Square anyway and watch:} $2x-1=9$, so $x=\blank{1.0cm}$. \\[3pt]
Check it: $\sqrt{2(5)-1}+5=\blank{1.0cm}+5=\blank{1.0cm}$, \\
which is not $2$. \; \textbf{Extraneous.} \\[3pt]
The algebra never lied --- it just never \emph{warned} you.
\end{minipage}
\end{center}
\end{notesbox}

\begin{notesbox}{4. Odd Index: The Trapdoor Closes}
Everything above was about \emph{even} roots. Change the index and the entire problem changes ---
which has been the theme of this unit since Lesson 6.0.

Cubing does not destroy sign information: every real number has exactly \blank{1.2cm} real cube
root, so $a^{3}=b^{3}$ forces $a=b$. Raising both sides to an \textbf{odd} power is
\blank{2.6cm}, and therefore it \textbf{cannot invent solutions}.
\begin{center}
\renewcommand{\arraystretch}{1.6}
\begin{tabularx}{0.96\linewidth}{@{}X|X@{}}
\textbf{$\sqrt{x-4}=-2$} & \textbf{$\sqrt[3]{x-4}=-2$} \\
\hline
A square root equal to a negative number: \blank{2.6cm} & Cube both sides: $x-4=\blank{1.4cm}$,
so $x=\blank{1.4cm}$ \\
& Check: $\sqrt[3]{\blank{1.4cm}}=\blank{1.2cm}$ \; \checkmark \\
\end{tabularx}
\end{center}
\emph{Same two numbers. Opposite outcomes.} The only difference is the \blank{1.4cm} --- which is
the sentence this whole unit has been built on.

\vspace{4pt}
\textbf{Worked example.} $2\sqrt[3]{x+1}-3=1$
\begin{center}
\renewcommand{\arraystretch}{1.7}
\begin{tabular}{@{}l l@{}}
$2\sqrt[3]{x+1}=\blank{1.2cm}$ & add $3$ \\
$\sqrt[3]{x+1}=\blank{1.2cm}$ & divide by $2$ --- \emph{now} the radical is alone \\
$x+1=\blank{1.2cm}$ & cube both sides \\
$x=\blank{1.2cm}$ & \\
\end{tabular}
\end{center}
\emph{Still check.} An odd index protects you from extraneous solutions, not from arithmetic
mistakes.
\end{notesbox}

\begin{notesbox}{5. Verifying on a Graph}
An equation like $\sqrt{x+4}-1=3$ is a question about Lesson 6.4's graphs: \emph{where does the
curve $y=\sqrt{x+4}-1$ meet the horizontal line $y=3$?}
\begin{center}
\begin{minipage}[c]{0.52\linewidth}
\centering
\begin{tikzpicture}[scale=0.40]
  \lgrid{-5}{15}{-3}{5}
  \begin{scope}
    \clip (-5,-3) rectangle (15,5);
    \draw[very thick, forest, domain=-4:15, samples=110, smooth] plot (\x,{sqrt((\x)+4)-1});
    \draw[very thick, navy] (-5,3)--(15,3);
    \draw[very thick, goldacc] (-5,-2)--(15,-2);
  \end{scope}
  \fill[forest] (-4,-1) circle (4pt);
  \fill[charcoal] (12,3) circle (3.4pt) node[above left, font=\scriptsize]{$(12,3)$};
  \node[navy, font=\scriptsize, right] at (12.6,3.7) {$y=3$};
  \node[goldacc, font=\scriptsize, right] at (12.6,-2.6) {$y=-2$};
\end{tikzpicture}
\end{minipage}
\begin{minipage}[c]{0.44\linewidth}
\small
$\sqrt{x+4}-1=3$ has \blank{1.2cm} solution, at $x=\blank{1.2cm}$ --- exactly what the algebra in
Section 1 gave. \\[6pt]
$\sqrt{x+4}-1=-2$ has \blank{1.2cm} solution(s), because the line $y=-2$ passes
\blank{2.4cm} the endpoint and the graph never reaches it. \\[6pt]
\textbf{The rule:} the number of \blank{2.6cm} between the two graphs is the number of real
solutions.
\end{minipage}
\end{center}
\vspace{2pt}
So if your algebra produces \emph{more} candidates than the picture has crossings, the extras are
\blank{2.6cm}. A graph will not do your algebra for you, but it will always tell you \emph{how
many} answers to expect.
\end{notesbox}

\begin{notesbox}{6. The Same Method in Exponent Clothing: $x^{m/n}=k$}
Lesson 6.1 said a radical is an exponent in disguise, so a radical equation is an
\emph{exponential-looking} equation in disguise too. To undo the power $\frac{m}{n}$, raise both
sides to its \textbf{reciprocal}, $\frac{n}{m}$ --- because
$\left(x^{m/n}\right)^{n/m}=x^{\,\blank{0.8cm}}=x$.
\begin{center}
\begin{minipage}[t]{0.46\linewidth}
\small
\textbf{(a)} \; $x^{3/2}=27$
\[ x=27^{\,\blank{0.9cm}}=\left(\sqrt[3]{27}\,\right)^{2}=\blank{1.2cm} \]
Check: $9^{3/2}=\left(\sqrt{9}\,\right)^{3}=\blank{1.2cm}$ \; \checkmark
\end{minipage}
\hfill
\begin{minipage}[t]{0.50\linewidth}
\small
\textbf{(b)} \; $x^{2/3}=9$
\[ x=\pm\,9^{\,\blank{0.9cm}}=\pm\left(\sqrt{9}\,\right)^{3}=\blank{1.8cm} \]
Both check! $(-27)^{2/3}=\left(\sqrt[3]{-27}\,\right)^{2}=(-3)^{2}=\blank{1.0cm}$
\end{minipage}
\end{center}
\vspace{3pt}
\begin{center}
\begin{tcolorbox}[colback=goldbg, colframe=goldacc, boxrule=0.7pt, arc=1.5mm, width=0.94\linewidth,
    left=3mm, right=3mm, top=1.5mm, bottom=1.5mm]
\small Read the fraction before you solve. An \textbf{even numerator} is a square in disguise, so
write the $\pm$ yourself. An \textbf{even denominator} is an even root, so the base cannot be
negative --- and a negative answer there would be \blank{2.6cm}.
\end{tcolorbox}
\end{center}
\end{notesbox}

\begin{practicebox}
Solve each equation. Show the check for every candidate, and state the solution set.

\vspace{4pt}
\textbf{1.} $\sqrt{3x-2}+4=9$
\vspace{0.55in}

\textbf{Solution set:} \blank{3.0cm}

\vspace{6pt}
\textbf{2.} $\sqrt{5x+6}=x$
\vspace{0.85in}

Candidates: \blank{3.0cm} \quad Extraneous: \blank{2.4cm} \quad \textbf{Solution set:}
\blank{2.4cm}

\vspace{4pt}
Why was that candidate extraneous? Say it in one sentence about \emph{signs}. \writeline
\end{practicebox}

\end{document}
